\section{Keycloak}

Um die Sicherheit der Authentifizierungs- und Autorisierungsanforderungen der Anwendung zu erfüllen, 
wurde Keycloak als Identity- und Access-Management-Lösung integriert. 
Keycloak bietet eine umfassende Plattform für die Verwaltung von Benutzern, 
Rollen und Berechtigungen sowie die Unterstützung verschiedener Authentifizierungsprotokolle 
wie OpenID Connect, OAuth 2.0 und SAML.



% ============================================================
% KEYCLOAK – GEGENSEITE DES JWT
% Empfehlung: In quarkus.tex oder architecture.tex einbauen
% ============================================================

\section{Keycloak und JWT-Authentifizierung}
\setauthor{Markus Tran}

Keycloak ist ein quelloffenes Identity- und Access-Management-System, das
Authentifizierung und Autorisierung als eigenständigen Dienst bereitstellt.
In der vorliegenden Anwendung übernimmt Keycloak die gesamte Verwaltung von
Benutzeridentitäten, sodass weder im Frontend noch im Backend eine eigene
Passwort- oder Benutzerverwaltung implementiert werden musste.

\subsection{Ablauf der Authentifizierung}

Der Authentifizierungsablauf folgt dem OAuth2-Protokoll mit dem Authorization
Code Flow. Wenn ein Benutzer auf einen geschützten Bereich der Anwendung
zugreift, leitet das Angular-Frontend ihn an Keycloak weiter. Nach erfolgreicher
Anmeldung stellt Keycloak ein signiertes JSON Web Token (JWT) aus, das der
Benutzer bei jeder weiteren Anfrage im \texttt{Authorization}-Header als
Bearer-Token mitsendet. Dieser Ablauf wird von \textit{Cufurovic} aus
Frontend-Sicht beschrieben, während der vorliegende Abschnitt die Verarbeitung
dieses Tokens im Backend behandelt.

\subsection{JWT-Verarbeitung im Backend}

Im Backend wird das vom Frontend mitgesendete JWT-Token von Quarkus mithilfe
der MicroProfile JWT-Spezifikation automatisch validiert und dekodiert. Ein
gültiges Token enthält unter anderem den sogenannten \textit{Subject Claim}
(\texttt{sub}), der die eindeutige Keycloak-UUID des authentifizierten Benutzers
enthält. Diese UUID wird im \texttt{UserService} über die injizierte
\texttt{JsonWebToken}-Instanz extrahiert:

\begin{lstlisting}[language=java, caption={Extraktion der Benutzer-UUID aus dem JWT im UserService}]
@Inject
JsonWebToken jwt;

public UUID getCurrentUserId() {
    return UUID.fromString(jwt.getSubject());
}
\end{lstlisting}

Die extrahierte UUID dient als primärer Schlüssel des Benutzers in der
Datenbank und stellt die Verbindung zwischen dem Keycloak-Identitätssystem
und dem internen Datenbankmodell her. Da die UUID direkt aus dem signierten
JWT stammt, ist eine Manipulation durch den Client ausgeschlossen.

\subsection{Rollenbasierte Zugriffskontrolle}

Neben der reinen Authentifizierung unterstützt Keycloak auch eine
rollenbasierte Zugriffskontrolle (Role-Based Access Control, RBAC). Im
Backend werden Endpunkte über die JAX-RS-Annotation \texttt{@RolesAllowed}
auf bestimmte Rollen beschränkt. Endpunkte, die benutzerspezifische Daten
liefern, sind auf die Rollen \texttt{user} und \texttt{medien-abo-users-role}
beschränkt, während öffentliche Endpunkte mit \texttt{@PermitAll} für alle
Aufrufer zugänglich sind. Das Frontend prüft denselben Rollenstatus, um
geschützte Navigationsbereiche ein- oder auszublenden -- beide Schichten
verwenden damit dieselbe Rolleninformation aus dem JWT.